% GENERATED FILE. Edit canonical lessons, then run npm run generate:books.
% canonical-source-hash: fnv1a64:14b36385594a38ef
% canonical-lessons: GE-C499-nadel, GE-C499-faden, GE-C499-stricken, GE-C499-haekeln, GE-C499-sticken

\chapter{Stricken und sticken}
\label{ch:ge-stricken-und-sticken}
\hlchaptermodality{499}

\begin{chapteropening}
\textbf{By the end of this chapter:} \emph{I can talk about knitting and needlework.}

Complete the last lesson of chapter 499: I can talk about knitting and needlework.
\end{chapteropening}

\section[die Nadel]{die Nadel — needle}

\label{lesson:GE-C499-nadel}



\begin{quote}
Before the new one: say the German for neighbourhood, then the German for hospitality.
\end{quote}

\subsection*{You'll want to know: die Nadel}
\textbf{die Nadel} — \textquotedblleft{}needle\textquotedblright{}.

\begin{quote}
Die Nadel ist so klein, dass meine Oma eine Brille braucht.
\end{quote}

\begin{quote}
\emph{The needle is so small that my grandma needs glasses.}
\end{quote}

Plural: \emph{die Nadeln}. Also a pin, or the needle of a pine tree.

\begin{cousinweb}
\emph{Nadel} is inherited from Middle High German \emph{nādele} and Old High German \emph{nādala}, from Proto-Germanic \emph{nēþlō} \textquotedblleft{}needle\textquotedblright{} (EtymDB); English \emph{needle} is a cousin.
\end{cousinweb}

\begin{grammarlens}[title={What you've built}]
One more word for telling, explaining and describing at B1.
\end{grammarlens}

\subsection*{Your turn}
\begin{itemize}
\raggedright
  \item \emph{Say it:} \emph{die Nadel}
  \item \emph{Say it:} \emph{die Nadel}, once more
  \item \emph{Say it:} a sentence of your own with \emph{die Nadel}
  \item \emph{Recall:} say the German for neighbourhood, then the German for hospitality, then say \emph{die Nadel} again
\end{itemize}

\begin{tcolorbox}[breakable,colback=teal!4,colframe=teal!35!black,title={Before you move on}]
What does die Nadel mean? (\textquotedblleft{}Needle\textquotedblright{}.) Say it once more.
\end{tcolorbox}
\section[der Faden]{der Faden — thread}

\label{lesson:GE-C499-faden}



\begin{quote}
Before the new one: say the German for hospitality, then the German for needle.
\end{quote}

\subsection*{You'll want to know: der Faden}
\textbf{der Faden} — \textquotedblleft{}thread\textquotedblright{}.

\begin{quote}
Der Faden ist zu kurz für den Knopf.
\end{quote}

\begin{quote}
\emph{The thread is too short for the button.}
\end{quote}

Plural: \emph{die Fäden}. \emph{Den Faden verlieren} — to lose the thread.

\begin{cousinweb}
\emph{Faden} is usually traced to Old High German \emph{fadum} \textquotedblleft{}the span of outstretched arms\textquotedblright{}, a length of thread; English \emph{fathom} is a cousin. EtymDB lists the German word without a source.
\end{cousinweb}

\begin{grammarlens}[title={What you've built}]
One more word for telling, explaining and describing at B1.
\end{grammarlens}

\subsection*{Your turn}
\begin{itemize}
\raggedright
  \item \emph{Say it:} \emph{der Faden}
  \item \emph{Say it:} \emph{der Faden}, once more
  \item \emph{Say it:} a sentence of your own with \emph{der Faden}
  \item \emph{Recall:} say the German for hospitality, then the German for needle, then say \emph{der Faden} again
\end{itemize}

\begin{tcolorbox}[breakable,colback=teal!4,colframe=teal!35!black,title={Before you move on}]
What does der Faden mean? (\textquotedblleft{}Thread\textquotedblright{}.) Say it once more.
\end{tcolorbox}
\section[stricken]{stricken — to knit}

\label{lesson:GE-C499-stricken}



\begin{quote}
Before the new one: say the German for needle, then the German for thread.
\end{quote}

\subsection*{You'll want to know: stricken}
\textbf{stricken} — \textquotedblleft{}to knit\textquotedblright{}.

\begin{quote}
Im Winter strickt meine Mutter Schals für die ganze Familie.
\end{quote}

\begin{quote}
\emph{In winter my mother knits scarves for the whole family.}
\end{quote}

Regular: \emph{hat gestrickt}.

\begin{cousinweb}
\emph{stricken} is built on Old High German \emph{stric} \textquotedblleft{}cord, loop\textquotedblright{} (EtymDB): knitting makes loops of yarn.
\end{cousinweb}

\begin{grammarlens}[title={What you've built}]
One more word for telling, explaining and describing at B1.
\end{grammarlens}

\subsection*{Your turn}
\begin{itemize}
\raggedright
  \item \emph{Say it:} \emph{stricken}
  \item \emph{Say it:} \emph{stricken}, once more
  \item \emph{Say it:} a sentence of your own with \emph{stricken}
  \item \emph{Recall:} say the German for needle, then the German for thread, then say \emph{stricken} again
\end{itemize}

\begin{tcolorbox}[breakable,colback=teal!4,colframe=teal!35!black,title={Before you move on}]
What does stricken mean? (\textquotedblleft{}To knit\textquotedblright{}.) Say it once more.
\end{tcolorbox}
\section[häkeln]{häkeln — to crochet}

\label{lesson:GE-C499-haekeln}



\begin{quote}
Before the new one: say the German for thread, then the German for to knit.
\end{quote}

\subsection*{You'll want to know: häkeln}
\textbf{häkeln} — \textquotedblleft{}to crochet\textquotedblright{}.

\begin{quote}
Sie häkelt eine kleine Mütze für das Baby.
\end{quote}

\begin{quote}
\emph{She is crocheting a little hat for the baby.}
\end{quote}

Done with one hooked needle, \emph{die Häkelnadel}.

\begin{cousinweb}
\emph{häkeln} is built on \emph{Häkel}, a small form of \emph{Haken} \textquotedblleft{}hook\textquotedblright{}; EtymDB lists the verb without a source.
\end{cousinweb}

\begin{grammarlens}[title={What you've built}]
One more word for telling, explaining and describing at B1.
\end{grammarlens}

\subsection*{Your turn}
\begin{itemize}
\raggedright
  \item \emph{Say it:} \emph{häkeln}
  \item \emph{Say it:} \emph{häkeln}, once more
  \item \emph{Say it:} a sentence of your own with \emph{häkeln}
  \item \emph{Recall:} say the German for thread, then the German for to knit, then say \emph{häkeln} again
\end{itemize}

\begin{tcolorbox}[breakable,colback=teal!4,colframe=teal!35!black,title={Before you move on}]
What does häkeln mean? (\textquotedblleft{}To crochet\textquotedblright{}.) Say it once more.
\end{tcolorbox}
\section[sticken]{sticken — to embroider}

\label{lesson:GE-C499-sticken}



\begin{quote}
Before the new one: say the German for to knit, then the German for to crochet.
\end{quote}

\subsection*{You'll want to know: sticken}
\textbf{sticken} — \textquotedblleft{}to embroider\textquotedblright{}.

\begin{quote}
Auf das Kissen hat sie eine Blume gestickt.
\end{quote}

\begin{quote}
\emph{She embroidered a flower on the cushion.}
\end{quote}

Regular: \emph{hat gestickt}. Not to be confused with \emph{stecken} — to put, to stick.

\begin{cousinweb}
\emph{sticken} is usually traced to Old High German \emph{sticken} \textquotedblleft{}to prick, to stitch\textquotedblright{}, related to \emph{stechen} \textquotedblleft{}to prick\textquotedblright{}; EtymDB lists the word without a source.
\end{cousinweb}

\begin{grammarlens}[title={What you've built}]
One more word for telling, explaining and describing at B1.
\end{grammarlens}

\subsection*{Your turn}
\begin{itemize}
\raggedright
  \item \emph{Say it:} \emph{sticken}
  \item \emph{Say it:} \emph{sticken}, once more
  \item \emph{Say it:} a sentence of your own with \emph{sticken}
  \item \emph{Recall:} say the German for to knit, then the German for to crochet, then say \emph{sticken} again
\end{itemize}

\begin{tcolorbox}[breakable,colback=teal!4,colframe=teal!35!black,title={Before you move on}]
What does sticken mean? (\textquotedblleft{}To embroider\textquotedblright{}.) Say it once more.
\end{tcolorbox}
